
%(BEGIN_QUESTION)
% Copyright 2016, Tony R. Kuphaldt, released under den Creative Commons Attribution License (v 1.0)
% This means you may do almost anything med denne work of mine, so long as you give me proper credit

Vedlikeholdet av et fungerende laboratorium er omfattende, spesielt i et program som instrumentering, der mye av utstyret kommer som donasjoner som må settes sammen, og der mange systemer er spesialbygde for formålet. Hver elev har ansvar for å bidra til vedlikehold av laboratoriet, fordi alle nyter godt av fasilitetene.

På siste dag i hvert kvartal settes det av tid til opprydding og reorganisering av laboratoriet. Dette er en full arbeidsdag med obligatorisk oppmøte. For å hjelpe elevene med å fokusere på oppgavene som må gjøres, viser listen nedenfor noe av arbeidet som er nødvendig for å klargjøre laboratoriet til neste kvartal. Oppgaver med tom linje foran blir tildelt lab-team for gjennomføring.

\vskip 10pt

\noindent
{\bf Lab tasks}

\begin{itemize}
\item{} Kontroller at alle små gjenstander med BTC-inventarmerke er malt i en tydelig farge, slik at de er enkle å finne ved årlig inventarsjekk.
\item{} \underbar{\hskip 50pt} Fei alle gulvflater i laboratoriet, og resirkuler eller kast avfallsmateriale.
\item{} \underbar{\hskip 50pt} Fei alle gulvflater i lagerrom, og resirkuler eller kast avfallsmateriale. Sett gjenstander som ligger på gulvet tilbake på riktig hylleplass.
\item{} \underbar{\hskip 50pt} Samle alle kobberrørbiter og legg dem i beholderen for kobber/messing-resirkulering.
\item{} \underbar{\hskip 50pt} Samle alt skrap av aluminium, rustfritt stål, messing og kobbertråd (biter kortere enn 1 fot) i metallskrap-bøttene ved nordvestlig utgang.
\item{} \underbar{\hskip 50pt} Lever resirkulerbart metall til lokal skraphandler, og ta med kontanter tilbake for å kjøpe pizza til lunsj.
\item{} \underbar{\hskip 50pt} Organiser oppbevaringskasser for fareskilt og maskeringstape. Samle inn ubrukte fareskilt fra laboratoriet og legg dem i riktig kasse.
\item{} \underbar{\hskip 50pt} Hjelp til med å lete etter manglende gjenstander i Team Tool Locker.
\item{} \underbar{\hskip 50pt} Rengjør alle arbeidsbenk- og bordflater.
\item{} \underbar{\hskip 50pt} Fjern gjenstander fra kompressorrommet, fei gulvet, og sørg for at rommet ikke brukes til tilfeldig lagring.
\item{} \underbar{\hskip 50pt} Samle kabellengder over 1 fot og legg dem i oppbevaringskasser inne i DCS-skapene for senere bruk.
\item{} \underbar{\hskip 50pt} Reorganiser området for kabelspoler: fjern tomme spoler fra stativet, og sørg for at bokser og løse spoler er ryddig stablet på gulvet.
\item{} \underbar{\hskip 50pt} Samle alle plastslanger og returner dem til riktig oppbevaringskasse.
\item{} \underbar{\hskip 50pt} Reorganiser skuffene for rørkoblinger (nordvesthjørnet i laboratoriet), slik at ingen rørgjengefittings er blandet inn, at alle fittings ligger i riktig skuff, og at alle skuffer er tydelig merket (med eksempelkoblinger på frontene).
\item{} \underbar{\hskip 50pt} Reorganiser skuffene for rørgjengefittings (nordvesthjørnet i laboratoriet), og sørg for at ingen rørkoblinger for slange er blandet inn.
\item{} \underbar{\hskip 50pt} Reorganiser skuffene for slangekoblinger (nordvesthjørnet i laboratoriet).
\item{} \underbar{\hskip 50pt} Reorganiser skuffene for rekkeklemmer og ice-cube-reléer (nordenden av laboratoriet).
\item{} \underbar{\hskip 50pt} Tapp kondensvann fra trykklufttank i kompressorrommet.
\item{} \underbar{\hskip 50pt} Sett alle bøker og manualer tilbake i bokhyllene.
\item{} \underbar{\hskip 50pt} Inspiser hvert reguleringspanel i laboratoriet, og fjern all kabling unntatt det som skal være permanent installert (120 VAC effekt, signalkabler mellom koblingsbokser). Sørg for at effektkabler til hver koblingsboks er godt festet og jordet.
\item{} \underbar{\hskip 50pt} Inspiser hver signalkoblingsboks i laboratoriet, og fjern all kabling unntatt det som skal være permanent installert (f.eks. 120 VAC effekt, signalkabler mellom koblingsbokser). Sørg for at effektkabler til hver koblingsboks er godt festet og jordet.
\item{} \underbar{\hskip 50pt} Kontroller merkingen på alle koblingsbokser og reguleringspaneler. Lag nye merker der gamle mangler, er skadet eller er vanskelige å lese.
\item{} \underbar{\hskip 50pt} Kontroller merkingen på alle permanent installerte kabler (f.eks. mellom koblingsbokser). Lag nye merker der gamle mangler, er skadet eller er vanskelige å lese.
\item{} \underbar{\hskip 50pt} Kontroller merkingen på alle rekkeklemmer inne i reguleringspaneler og koblingsbokser. Lag nye merker der gamle mangler eller er vanskelige å lese.
\item{} \underbar{\hskip 50pt} Fjern alt avfall inne i reguleringspaneler og koblingsbokser i hele laboratoriet, bruk støvsuger ved behov.
\item{} \underbar{\hskip 50pt} Rengjør deadweight-testere (de lekker ofte olje). {\it Hint: WD-40 fungerer godt som løsemiddel for å fjerne oljesøl.}
\item{} \underbar{\hskip 50pt} Utfør vedlikehold på turbokompressorsystemet: sikker utkobling/tagging, kontroll av oljenivå, reparasjon av oljelekkasjer, utbedring av dårlige kabeltilkoblinger, rengjøring i reguleringsskap, og etterstramming av effektterminaler.
\item{} \underbar{\hskip 50pt} Returner alt fellesverktøy (f.eks. effektbor og sager) til riktig lagringsplass (håndverktøy i verktøyskuffen i nordøsthjørnet av laboratoriet, effektverktøy på verktøyhyllen i lagerområdet oppe).
\item{} \underbar{\hskip 50pt} Fjern gjenstander fra alle lagerskap i nordenden av laboratoriet, rydd hyller for skrot (f.eks. pH-prober som har blitt liggende tørre), og sett alt tilbake på riktig plass. Monter deksler på transmittere som mangler dette, spesielt pneumatiske transmittere som er sårbare uten deksel.
\item{} \underbar{\hskip 50pt} Utfør visuell kontroll av alle allmenntrykkregulatorer i nordlige lagringshyller for manglende justeringsbolter, manglende slangekoblinger, skadede gjenger osv. Reparer ved behov.
\item{} \underbar{\hskip 50pt} Test alle trykktransmittere som ikke er merket ``good'' for å bekrefte om de faktisk er defekte. Reparer om mulig, berg deler og kasser hvis ikke. {\it Ikke kast instrumenter med BTC-inventarmerke!}
\item{} \underbar{\hskip 50pt} Test alle temperaturtransmittere som ikke er merket ``good'' for å bekrefte om de faktisk er defekte. Reparer om mulig, berg deler og kasser hvis ikke. {\it Ikke kast instrumenter med BTC-inventarmerke!}
\item{} \underbar{\hskip 50pt} Test alle I/P-omformere som ikke er merket ``good'' for å bekrefte om de faktisk er defekte. Reparer om mulig, berg deler og kasser hvis ikke. {\it Ikke kast instrumenter med BTC-inventarmerke!}
\item{} \underbar{\hskip 50pt} Test alle presisjonstrykkmålere som ikke er merket ``good'' for å bekrefte om de faktisk er defekte. Reparer om mulig, berg deler og kasser hvis ikke. {\it Ikke kast instrumenter med BTC-inventarmerke!}
\item{} \underbar{\hskip 50pt} Test alle presisjonstrykkregulatorer som ikke er merket ``good'' for å bekrefte om de faktisk er defekte. Reparer om mulig, berg deler og kasser hvis ikke. {\it Ikke kast instrumenter med BTC-inventarmerke!}
\item{} \underbar{\hskip 50pt} Returner alle feltinstrumenter (f.eks. transmittere) og diverse enheter (f.eks. trykkmålere og regulatorer) til riktig lagringsplass. {\it Merk at I/P-transdusere og ventilposisjonerere skal bli stående ved sine respektive reguleringsventiler, ikke legges på lager.}
\item{} \underbar{\hskip 50pt} Lagre alle 2$\times$2 fot kryssfiner-prosesstavler på sikre steder, og sørg for at hver tavle er klar til bruk neste kvartal.
\item{} \underbar{\hskip 50pt} Sørg for at hver reguleringsventil montert på stativene i laboratoriet har en I/P-transduser montert i nærheten, med Swagelok-koblinger i god stand for tilkobling av trykkluft og signal.
\item{} \underbar{\hskip 50pt} Kontroller at hver ventil er forsvarlig festet til stativet, og at eventuell posisjoner har tilbakekoblingsarm korrekt koblet til ventilspindel (ingen manglende fjærer, bøyde stag eller tydelig feilstilling).
\item{} \underbar{\hskip 50pt} Fjern innhold fra skapet for brannfarlige væsker, tørk av hyller for væske og rester, og fyll opp igjen på en ryddig og sikker måte.
\item{} \underbar{\hskip 50pt} Rengjør alle grillene for rester etter lunsjer og innsamlingsarrangement. {\it Merk: det kan være nødvendig å ta med grillrister og fettoppsamlingsbrett til bilvask og bruke avfettingsmiddel for å få dem tilstrekkelig rene.}
\item{} \underbar{\hskip 50pt} Tilbakestill alle funksjonsblokkparametere i DCS ``Generic Loops'' til standardverdier. Se dokumentasjonen på hovedmaskinen BTC\_PPlus for riktige parameterverdier.
\item{} \underbar{\hskip 50pt} Kontroller manometrene på kalibreringsbenken: de med rød væske skal ha korrekt nivå, og alle øvrige (normalt fylt med destillert vann) skal være helt tømt.
\item{} \underbar{\hskip 50pt} Slå på trykkluft til kalibreringsbenken, kontroller lekkasjer, og verifiser at alle trykkregulatorer fungerer korrekt.
\item{} \underbar{\hskip 50pt} Rengjør kjøleskap og kast matvarer som fortsatt står igjen.
\item{} \underbar{\hskip 50pt} Rengjør alle stekeovner og annet kjøkkenutstyr grundig.
\item{} \underbar{\hskip 50pt} Sett alle hylleplater tilbake på riktig plass i stativene.
\item{} \underbar{\hskip 50pt} Rengjør og reorganiser alle hyller i klasserom DMC 130 med komponenter for praktiske mastery-vurderinger. Kast skadde jumperledninger, batteriklemmer osv. Kast batterier der terminalspenningen er under 80\% av merkespenningsverdien (f.eks. under 7.2 V for et 9 V-batteri).
\item{} \underbar{\hskip 50pt} Koble ut effekt til alle reguleringssystemer unntatt DCS.
\item{} \underbar{\hskip 50pt} Plasser donerte komponenter på riktig lagringsplass.
\item{} \underbar{\hskip 50pt} Rengjør alle whiteboard-tavler med egnet rengjøringsmiddel slik at de faktisk blir hvite igjen.
\end{itemize}

\filbreak

\begin{itemize}
\item{} {\it Lærere kan legge til punkter i listen ved behov:}
\vskip 10pt
\item{} \underbar{\hskip 50pt} 
\vskip 10pt
\item{} \underbar{\hskip 50pt} 
\vskip 10pt
\item{} \underbar{\hskip 50pt} 
\vskip 10pt
\item{} \underbar{\hskip 50pt} 
\vskip 10pt
\item{} \underbar{\hskip 50pt} 
\vskip 10pt
\item{} \underbar{\hskip 50pt} 
\vskip 10pt
\item{} \underbar{\hskip 50pt} 
\end{itemize}


\filbreak

\noindent
{\bf Personal tasks}

\begin{itemize}
\item{} Før opp ``syketimer'' for tapt tid dette kvartalet (husk at dette {\it ikke} gjøres automatisk for deg!).
\item{} Doner ubrukte ``syketimer'' til medstudenter som trenger dem.
\item{} Ta igjen eventuelle quizer du har gått glipp av på grunn av fravær dette kvartalet (husk at en quiz som ikke tas, regnes som stryk).
\end{itemize}

\underbar{file i01229}
%(END_QUESTION)





%(BEGIN_ANSWER)

 
%(END_ANSWER)





%(BEGIN_NOTES)

Jeg anbefaler å forhåndstildele minst to oppgaver til hvert lab-team, slik at ikke noen blir stående uten å gjøre noe mens andre gjør alt arbeidet.

%INDEX% Lab-opprydding og reorganisering, arbeidsliste

%(END_NOTES)
